\begin{examplebox}

\textbf{\textcolor{CExample}{例 1（基础）}}

\textbf{\textcolor{CExample}{题目：}} 已知两个球的半径之比为 $2:3$，求它们的表面积之比和体积之比。

\textbf{\textcolor{CExample}{解题步骤：}}
\begin{description}[style=nextline, leftmargin=2.8em, labelsep=0.5em, itemsep=0.45em]
    \item[\textbf{第一步（识别比例）：}] 半径比 $m:n = 2:3$。
    \par\textbf{理由：} 直接提取输入条件。
    \item[\textbf{第二步（套用面积比公式）：}] 表面积之比
    \[
    S_1:S_2 = m^2:n^2 = 2^2:3^2 = 4:9.
    \]
    \par\textbf{理由：} 表面积比等于半径比的平方。
    \item[\textbf{第三步（套用体积比公式）：}] 体积之比
    \[
    V_1:V_2 = m^3:n^3 = 2^3:3^3 = 8:27.
    \]
    \par\textbf{理由：} 体积比等于半径比的立方。
\end{description}

\textbf{\textcolor{CExample}{关键结论：}} \textbf{表面积比 $4:9$，体积比 $8:27$。}

\medskip
\textbf{\textcolor{CExample}{例 2（稍变形）}}

\textbf{\textcolor{CExample}{题目：}} 两个球的半径分别为 $3\,\text{cm}$ 和 $5\,\text{cm}$，求大球体积是小球体积的多少倍，并求大球体积比小球体积大多少倍。

\textbf{\textcolor{CExample}{解题步骤：}}
\begin{description}[style=nextline, leftmargin=2.8em, labelsep=0.5em, itemsep=0.45em]
    \item[\textbf{第一步（明确顺序）：}] 小球半径为 $3\,\text{cm}$，大球半径为 $5\,\text{cm}$，所以
    \[
    R_{\text{小}}:R_{\text{大}}=3:5.
    \]
    \par\textbf{理由：} 先把“小球”和“大球”的顺序说清楚，避免比例方向写反。

    \item[\textbf{第二步（套用体积比公式）：}] 小球与大球的体积比为
    \[
    V_{\text{小}}:V_{\text{大}}=3^3:5^3=27:125.
    \]
    因此
    \[
    \frac{V_{\text{大}}}{V_{\text{小}}}=\frac{125}{27}.
    \]
    \par\textbf{理由：} 体积比等于半径比的立方，同时注意题目问的是“大球相对于小球”。

    \item[\textbf{第三步（区分“是几倍”和“大几倍”）：}] 大球体积是小球体积的
    \[
    \frac{125}{27}\approx 4.63
    \]
    倍；大球体积比小球体积大
    \[
    \frac{V_{\text{大}}-V_{\text{小}}}{V_{\text{小}}}
    =\frac{125-27}{27}
    =\frac{98}{27}
    \]
    倍。
    \par\textbf{理由：} “A 是 B 的多少倍”用 $\dfrac{A}{B}$；“A 比 B 大多少倍”通常用 $\dfrac{A-B}{B}$。考试中若题干表达不严谨，优先看题目给出的口径或老师要求。
\end{description}

\textbf{\textcolor{CExample}{关键结论：}} \textbf{大球体积是小球体积的 $\dfrac{125}{27}$ 倍；大球体积比小球体积大 $\dfrac{98}{27}$ 倍。}

\medskip
\textbf{\textcolor{CExample}{例 3（综合应用）}}

\textbf{\textcolor{CExample}{题目：}} 两个球的半径之比为 $1:2$，它们的体积之和为 $36\pi$，求较大球的表面积。

\textbf{\textcolor{CExample}{解题步骤：}}
\begin{description}[style=nextline, leftmargin=2.8em, labelsep=0.5em, itemsep=0.45em]
    \item[\textbf{第一步（设半径与体积比）：}] 设小球半径为 $r$，大球半径为 $2r$。则体积比
    \[
    V_{\text{小}}:V_{\text{大}}=1^3:2^3=1:8.
    \]
    \par\textbf{理由：} 体积比等于半径比的立方。

    \item[\textbf{第二步（按体积比分配体积和）：}] 设小球体积为 $V$，则大球体积为 $8V$。由题意
    \[
    V+8V=9V=36\pi,
    \]
    解得 $V=4\pi$，所以大球体积为
    \[
    V_{\text{大}}=8V=32\pi.
    \]
    \par\textbf{理由：} 体积和已知，先按 $1:8$ 分配。

    \item[\textbf{第三步（求大球半径）：}] 设大球半径为 $R$。由
    \[
    \frac{4}{3}\pi R^3=32\pi
    \]
    得
    \[
    R^3=24=8\times 3,
    \qquad
    R=2\sqrt[3]{3}.
    \]
    \par\textbf{理由：} 用大球体积直接反求大球半径，减少中间变量。

    \item[\textbf{第四步（求大球表面积）：}]
    \[
    S_{\text{大}}=4\pi R^2
    =4\pi\left(2\sqrt[3]{3}\right)^2
    =16\pi\cdot 3^{2/3}.
    \]
    \par\textbf{理由：} 代入球表面积公式 $S=4\pi R^2$。
\end{description}

\textbf{\textcolor{CExample}{关键结论：}} \textbf{较大球的表面积为 $16\pi\cdot 3^{2/3}$。}

\end{examplebox}
